\section{Introduction}
The general matrix-valued transport problem of Li and Zou combines a
constant-coefficient continuity operator with a weighted quadratic transport
and reaction action \cite{LiZou}. It includes the Kantorovich--Bures model of
Brenier and Vorotnikov \cite{BrenierVorotnikov} and the spatial matricial
interpolation model of Chen, Georgiou and Tannenbaum \cite{ChenGeorgiouTannenbaum}.
The choice of operator can affect limiting costs and convergence behavior,
as the examples below show.

The paragraph following Proposition 4.6 of \cite{LiZou} identifies vanishing
reaction weights, large transport weights, and partial eigenvalue degeneration
as further limiting regimes. We retain their signed-Radon support formulation
and finite positive semidefinite endpoints, extending the measure-level
parameter analysis to singular reaction weights as well.

The weights enter the action through their Gram matrices. Under an eventual
relative comparison with factors tending to one, Theorem~\ref{thm:finite}
gives stability of extended minimum values for simultaneously convergent
weights. For admissible paths with uniformly bounded actions, bounded weights give
coordinate-measure compactness; a small scalar shift of a continuous mass-dual
field proves the joint lower bound; relative comparison supplies constant-path
recovery. The comparison is a sufficient condition; fixed eigenbases give it
automatically. The reaction-zero
regime is a corollary, including an infinite balanced value.

For diverging transport weights, we preserve the finite-action class and prove a direct
infimum limit. Theorems~\ref{thm:ray} and~\ref{thm:uniform} treat fixed
singular rays and uniform divergence in all directions separately.
For entrywise divergence on a convex domain with a
positive-definite reaction weight, spatial aggregation and a three-part hub
construction identify the uniform-divergence limit explicitly
(Theorem~\ref{thm:hub}). The Bures formula,
its Procrustes representation \cite{BhatiaJainLim}, and weighted reaction
congruence \cite[Remark 3.11]{LiZou} are existing tools; related weighted
Bures geometry is studied in \cite{HanEtAl}.

The examples in Section~\ref{sec:examples} illustrate the scope of these
results. Rotation can create
a finite excess cost or divergence in a partially degenerating reaction
problem, and a jump in a scalar transport problem. Divergence of the Frobenius
norm does not imply convergence of minimum values. Even isotropic divergence
need not produce a total-matrix Bures value for the Kantorovich--Bures operator.
Finally, Section~\ref{sec:separation} gives a high-frequency family of smooth,
uniformly positive densities that separates convergence in two concrete
unit-weight models.

Scalar dynamic transport/reaction limits already appear in \cite{ChizatEtAl}.
The arguments use established convex duality, measure compactness and
finite-dimensional matrix geometry. The discrete convergence studied in
\cite{LiZouDiscrete} concerns a different limiting parameter, whereas the
present work varies the transport and reaction weights of the continuous
problem.
